Ein schwimmender Eiswürfel der Masse $m$ verdrängt so viel Flüssigkeit, dass deren Gewichtskraft gleich seiner eigenen ist, also eine Flüssigkeitsmasse $m$. Im Wasser ist das verdrängte Volumen
\begin{align}
 V_\text{verdr} &= \frac{m}{\rho_\text{W}}
\end{align}
Das Schmelzwasser hat ebenfalls die Masse $m$ und die Dichte $\rho_\text{W}$, also das Volumen $V_\text{S} = m/\rho_\text{W} = V_\text{verdr}$. Es füllt genau den Raum, den das Eis unter Wasser eingenommen hat: Der Wasserspiegel \result{\text{bleibt gleich}}.

Ein stark alkoholhaltiges Getränk hat eine kleinere Dichte $\rho_\text{D} < \rho_\text{W}$. Schwimmt das Eis noch ($\rho_\text{D} > \rho_\text{Eis}$), verdrängt es das Volumen $m/\rho_\text{D}$; ist das Getränk sogar weniger dicht als Eis, sinkt der Würfel und verdrängt sein ganzes Volumen $m/\rho_\text{Eis}$. In beiden Fällen ist das verdrängte Volumen grösser als das Volumen $m/\rho_\text{W}$ des Schmelzwassers. Der Flüssigkeitsspiegel \result{\text{sinkt}} beim Schmelzen. (Beim Mischen von Wasser und Alkohol nimmt das Volumen zudem etwas ab, was den Effekt verstärkt.)
